
\subsection{Protocol Motivation}

\textbf{UltraGroth} is addressed to insignificantly modify Groth16 protocol to
enable sampling random field elements effectively and further pass them as
signals to the circuit. The primary implication of such possibility is, of
course, lookup checks!

Moreover, such modification saves all the verification benefits of Groth16: in
UltraGroth's basic form, it adds only a single $\mathbb{G}_1$ point to the
proof, a single pairing operation, and a single hashing operation per
randomness. 

One of the examples where UltraGroth has already been used is
\href{https://mirror.xyz/0x90699B5A52BccbdFe73d5c9F3d039a33fb2D1AF6/T4lHiBlo7VoYp-5SEuitQnY7_ullpwkZbanOOwjeZI4}{\textit{Bionetta}}
-- an R1CS-based zkML framework used by
\href{https://rarimo.com/}{\textit{Rarimo}}. When implementing
$\text{ReLU}(x)=\max\{0,x\}$ non-linearity in-circuit, in the vanilla Groth16
one needs to perform the full bit-decomposition of $x$, thus costing $\approx
254$ constraints per call. As a result, the circuit size for the neural network
is roughly $\mathcal{O}(n)$ with $n=254L$ with $L$ being the number of ReLU
calls. Using UltraGroth, we reduced the number of constraints to the
\textit{sublinear complexity}: $\mathcal{O}(n/\log n)$.

\subsubsection{Historical Notes}

Before delving into specifics, we briefly mention the history of this protocol,
which is quite curious. This blog is based on the original
\href{https://hackmd.io/@Merlin404/Hy_O2Gi-h}{Lev Soukhanov's post} in 2023.
However, as we prepared the conference paper for Bionetta, we discovered that
his construction is nothing but a generalization of
the $\textsf{MIRAGE}$ protocol \cite{mirage}, posted back
in 2020! (in subsequent notation, $\textsf{MIRAGE}$ is an UltraGroth protocol
with $d=1$). Even more surprisingly, it seems that independently of Lev
Soukhanov's construction, Alex Ozdemir, Evan Laufer, and Dan Boneh introduced
the $\textsf{MIRAGE+}$ protocol \cite{mirageplus} in 2024,
which was used for proving RAM computations correctness. This protocol, which is
hard to believe, but also coincides with the Lev Soukharnov's construction
(although I believe the former generalizes the construction a bit more
rigorously). That said, to the best of my knowledge, Lev Soukhanov's
construction is the first to generalize $\textsf{MIRAGE}$ to the multiple-round
interactive protocol.

All three papers (Bionetta included) proved completeness, soundness, and
zero-knowledge of this construction, so the protocol that follows can be safely
integrated into production systems. We, of course, drop formalities in this blog
and recommend checking \cite{mirageplus} for proof specifics in case someone is interested.

Finally, our personal opinion is that this construction is vastly underestimated
and it is surprising that it is not yet used in production. Indeed, integrating
lookup checks comes with a tiny cost of additional pairing which enables
wrapping non-native protocols in Groth16 (e.g., zk-STARKs) \textit{much} more
effectively (up to a factor of $\times 20$). 

\begin{remark}
The subsequent explanations will be quite technical, so if you want easier and
more straightforward overview, you can additionally check out
\href{https://github.com/ZamDimon/presentations-and-talks/blob/main/ultragroth/ultragroth-static.pdf}{\textit{our
brief internal Distributed Lab presentation}}.
\end{remark}

\subsection{Lookup Checks}

Here we restate the key fact that we explored in \Cref{section:lookups}.

\begin{theorem} 
    Given two sequences of elements $\{t_i\}_{i \in [v]}$ and $\{z_i\}_{i \in
    [n]}$, the inclusion check $\{z_i\}_{i \in [n]} \subseteq \{t_i\}_{i \in
    [v]}$ is satisfied with overwhelming probability if and only if there exists
    the set of multiplicities $\{\mu_i\}_{i \in [v]}$ where $\mu_i = \#\{j \in
    [n]: z_j = t_i\}$ such that for randomly sampled $\gamma \leftarrowS
    \mathbb{F}$:
\[
    \sum_{i \in [n]} \frac{1}{\gamma+z_i} = \sum_{i \in [v]} \frac{\mu_i}{\gamma+t_i}
\]
\end{theorem}

More specifically, checking such equality at a random point from $\mathbb{F}$
results in the soundness error of up to $(n+v)/|\mathbb{F}|$, which becomes
negligible for large $|\mathbb{F}|$.

Now, the right hand size of the equation, given that we have sampled the
randomness $\gamma \gets \mathbb{F}$, can be computed in $v$ constraints, while
the left hand side is computed in additional $n$ constraints (recall that
inversion costs a single constraint to prove). Again, since Groth16 is a
non-interactive protocol in its essence (meaning, compared to PlonK or
sumcheck-based protocols, it is not constructed by converting an interactive
protocol to non-interactive), there is no yet effective way to sample $\gamma$
(besides probably in-circuit hashing which is insanely costly). Yet, we do
propose a way to sample $\gamma$ by turning the Groth16 into the
"interactive-like" proof system.

The proceeding \textit{UltraGroth} construction is essentially devoted to
sampling $\gamma$ and passing it as a signal to the circuit soundly.

\subsubsection{Some other implications}

Sampling randomnesses overall enables numerous applications besides lookup
checks. For instance, suppose you have matrices $A,B \in \mathbb{F}^{n \times
n}$ and you want to write the circuit that computes their product: $C=AB$. The
naive way to do that is to compute it by definition. However, such circuit costs
$\mathcal{O}(n^3)$ constraints, which is a lot. Turns out that sampling
randomnesses can reduce this significantly. Indeed, one can apply the
\textit{Freiveld's} protocol: first sample $\boldsymbol{\gamma} \leftarrowS \mathbb{F}^n$ and
instead of checking $AB=C$, check whether matrices $AB$ and $C$ project $\boldsymbol{\gamma}$
to the same vector (this looks somewhat similar to Schwartz-Zippel lemma, but in
the world of matrices):
\[
(AB)\boldsymbol{\gamma} = C\boldsymbol{\gamma} \iff A(B\boldsymbol{\gamma}) = C\boldsymbol{\gamma}
\]

Now notice that this can be implemented in approximately $3n^2$ constraints: one
first multiplies $u \gets B\boldsymbol{\gamma}$, then multiplies $Au$ and compares it to the
product $C\boldsymbol{\gamma}$. Note that matrix-vector product can be implemented in $n^2$
constraints, which is much more effective than matrix-matrix product with $n^3$
constraints.

However, currently, in Circom there is no way to pass $\boldsymbol{\gamma}$ as a signal to
the circuit. Indeed, one can write \texttt{signal input gamma[n]} and mark it as
public, but how can the verifier be ensured that corresponding values in witness
were indeed sampled uniformly randomly by the prover? So let's see how to fix
it.

\subsection{UltraGroth}

\subsection{Two-round Interactive R1CS}

Suppose we split our R1CS into multiple rounds. Consider the case where we want
to implement lookup checks. In such case, it would be nice if we could compile
the following interactive protocol (IP) into the zk-SNARK for a relation
$\mathcal{R}_{\text{QAP}}$:
\begin{itemize}
    \item \textit{Round 0}: The prover runs the circuit without imposing lookup
    checks. The prover sends resultant witness $\mathbbm{w}_0$ to verifier.
    Verifier responds with the random challenge $x_1 \leftarrowS \mathbb{F}$.
    \item \textit{Round 1}: The prover computes the witness $\mathbbm{w}_1$ that
    corresponds to the lookup check (that is, that verifies equality $\sum_{i}
    \frac{1}{x_1+z_i}=\sum_i \frac{\mu_i}{x_1+t_i}$) using the randomness $x_1$
    provided by verifier. The prover sends $\mathbbm{w}_1$ and polynomial $h(X)$
    according to $\mathcal{R}_{\text{QAP}}$.
    \item Finally, the verifier checks the same QAP equation
    $\ell(X)r(X)=o(X)+t_{\Omega}(X)h(X)$.
\end{itemize}

We would like to compile such IP into zk-SNARK by applying the Fiat-Shamir
heuristic. However, here is the issue: $x_1$ must be derived by hashing the
previous transcript, which is the public statement $\mathbbm{x}_0$ and witness
$\mathbbm{w}_0$. However, how can we hash the witness $\mathbbm{w}_0$, given that
the whole Groth16 proving system compiles the verification equation to just
three points?

Now, suppose we partition the index set of public signals $\mathcal{I}_X$ into
two parts (corresponding to each round): $\mathcal{I}_X^{\langle 0 \rangle}$ and
$\mathcal{I}_X^{\langle 1 \rangle}$. Corresponding private signals are indexed
by $\mathcal{I}_W^{\langle 0 \rangle}$ and $\mathcal{I}_W^{\langle 1 \rangle}$
which partition $\mathcal{I}_W$. In the example above, $\mathcal{I}_X^{\langle 0
\rangle}$ corresponds to the regular public signals in relation
$\mathcal{R}_{\text{QAP}}$ while $\mathcal{I}_X^{\langle 1 \rangle}$ indexes the
single solution witness element, corresponding to sampled randomness $x_1$.

The idea is the following: we split the polynomial $c(X)$ (which was described
in Groth16 construction above) into two pieces $c_0(X)$ and $c_1(X)$. The first
piece corresponds to the first round indexed by $(\mathcal{I}_X^{\langle 0
\rangle}, \mathcal{I}_W^{\langle 0 \rangle})$ and is formed as follows: sample
$r_0 \leftarrowS \mathbb{F}$, add toxic parameter $\delta_0$, and compute the point
$\pi_C^{\langle 0 \rangle} \gets g_1^{c_0(\tau)}$ where:
\[
c_0(X) = \delta_0^{-1}\sum_{j \in \mathcal{I}_W^{\langle 0 \rangle}} z_jf_j(X) + r_0\delta
\]
Notice that the point $\pi_C^{\langle 0 \rangle}$ encodes all witness elements
in Round 0. Thus, to sample randomness, we can hash prover parameters and this
point: $x_1 \gets \mathcal{H}(\mathsf{pp}, \pi_C^{\langle 0 \rangle})$. Then,
points $\pi_A=g_1^{a(\tau)}$ and $\pi_B=g_1^{b(\tau)}$ are computed in the same
way as in the regular Groth16, while the third piece $\pi_C^{\langle 1 \rangle}
= g_1^{c_1(\tau)}$ is computed using:
\[
c_1(X) = \delta^{-1}\left(\sum_{j \in \mathcal{I}_W^{\langle 1 \rangle}} z_jf_j(X)+h(X)t_{\Omega}(X)\right) + a(X)s + b(X)r - r_0\delta_0 - rs\delta
\]
Essentially, such construction ensures that $\delta_0c_0(X)+\delta c_1(X)$ gives
exactly the polynomial $\delta c(X)$ in the original Groth16. This allows the
prover to construct the proof as a tuple $\pi=(\pi_A,\pi_B,\pi_C^{\langle 0
\rangle},\pi_C^{\langle 1 \rangle})$ and verifier to check it using:
\[
e(\pi_A,\pi_B) = e(g_1^{\alpha},g_2^{\beta}) \cdot e(g_1^{i(\tau)},g_2^{\gamma}) \cdot e(\pi_C^{\langle 0 \rangle}, g_2^{\delta_0}) \cdot e(\pi_C^{\langle 1 \rangle}, g_2^{\delta}),
\]
where $i(X) = \gamma^{-1}\sum_{i \in \mathcal{I}_X}z_if_i(X)$ as before. As in
Groth16, the term $e(g_1^{\alpha}, g_2^{\beta})$ can be pre-computed, thus the
cost of verification is 4 pairing operations.

\begin{remark}
The aforementioned protocol, as mentioned earlier, is essentially the $\textsf{MIRAGE}$ protocol, and can be implemented as is. It is complete, sound, and zero-knowledge. In the subsequent section, we will generalize this protocol for more than 1 round.
\end{remark}

\subsection{Generalizing to multiple rounds}

To generalize the aforementioned protocol, we define the $(d+1)$-round quadratic arithmetic program (or $d$QAP, for short) as follows:
\[
\small\mathcal{R}_{\text{dQAP}} = \left\{ 
    \begin{matrix} 
        \mathbbm{x}_i = \{z_j\}_{j \in \mathcal{I}_{X}^{\langle i \rangle}} \\
        \mathbbm{w}_i = \{z_j\}_{j \in \mathcal{I}_{W}^{\langle i \rangle}} \\
        \text{for }i \in [d+1]
    \end{matrix}
    \; \middle| \; 
    \begin{matrix}
        \ell(X) \cdot r(X) = o(X) + t_{\Omega}(X)h(X) \\
        \ell(X) = \sum_{i \in [n]}z_i\ell_i(X), \\
        r(X) = \sum_{i \in [n]}z_ir_i(X), \\
        o(X) = \sum_{i \in [n]}z_io_i(X), \\
        \text{for some} \; h(X) \in \mathbb{F}[X]
    \end{matrix}
\right\},
\]
where indices $\{\mathcal{I}_X^{\langle i \rangle}\}_{i \in [d+1]}$ and $\{\mathcal{I}_W^{\langle i \rangle}\}_{i \in [d+1]}$ partition $[n]$. In this notation, all $\{\mathbbm{w}_i\}_{i \in [d+1]}$ are called \textit{private witnesses}, $\mathbbm{x}_0$ is a \textit{public statement}, while $\{\mathbbm{x}_i\}_{i \in [d+1]\setminus \{0\}}$ are \textit{challenges} (that are going to be sampled by means of UltraGroth).

Since prover and verifier participate in the interactive protocol, the prover cannot compute the whole witness on his own without verifier's randomnesses. For that reason, we define \textit{strategy} as the collection of functions $S=\{S_i\}_{i \in [d]}$, each of which computes the witness for the given round given previous witnesses and challenges and the current challenge, sampled by the verifier. In other words,
\[
\mathbbm{w}_i = S_i(\mathbbm{x}_0,\dots,\mathbbm{x}_i,\mathbbm{w}_0,\dots,\mathbbm{w}_{i-1})
\]

\begin{example}
For instance, 0QAP represents the regular Groth16 construction, where strategy
$S_0$ is a regular witness calculator: $\mathbbm{w}_0=S_0(\mathbbm{x}_0)$. The
protocol which involes lookup checks is a SNARK over 1QAP, where $S_0$ computes
the R1CS witness without lookup constraints, while $S_1$, based on the sampled
randomness(es) $\mathbbm{x}_1$, imposes lookup check condition:
$\mathbbm{w}_1=S_1(\mathbbm{x}_0,\mathbbm{x}_1,\mathbbm{w}_0)$.
\end{example}

Now, we can describe the \textit{UltraGroth} protocol over $\mathcal{R}_{\text{dQAP}}$.

\textbf{Setup($1^{\lambda},\mathcal{R}_{\text{dQAP}}$).} As in previous sections, denote by $f_i(X):= \beta \ell_i(X)+\alpha r_i(X)+o_i(X)$ the random linear combination of $\ell_i,r_i,o_i$. During the trusted setup, one generates the toxic waste $\alpha,\beta,\gamma,\{\delta_i\}_{i \in [d+1]},\tau \leftarrowS \mathbb{F}$ and computes the following common parameters:
\begin{align*}
\mathsf{pp} = \Big( &g_1^{\alpha},g_1^{\beta},\{g_1^{\delta_i}\}_{i \in [d]},g_2^{\beta}, g_2^{\gamma}, \{g_2^{\delta_i}\}_{i \in [d]}, \\
&\Big\{g_1^{\tau^i}, g_2^{\tau^i}, g_1^{\tau^it(\tau)/\delta_d}\Big\}_{i \in [n]}, \left\{g_1^{f_i(\tau)/\gamma}\right\}_{i \in \mathcal{I}_X}, \left\{g_1^{f_j(\tau)/\delta_i}\right\}_{(i,j) \in [d] \times \mathcal{I}_W^{(i)}}\Big)
\end{align*}
\[
\mathsf{vp} = \left( g_1,g_2, g_2^{\gamma}, \{g_2^{\delta_i}\}_{i \in [d]}, g_T^{\alpha\beta}, \left\{g_1^{f_i(\tau)/\gamma}\right\}_{i \in \mathcal{I}_{X}} \right)
\]

\textbf{Prove($\mathsf{pp},\mathbbm{x}_0,S$)}. The prover initializes the accumulator $a_0 := \mathcal{H}(\mathsf{pp})$. For each round $i \in [d]$ she conducts the following steps:
\begin{itemize}
    \item Samples $r_i \leftarrowS \mathbb{F}$
    \item Computes witness $\mathbbm{w}_i \gets S_i(\mathbbm{x}_0,\dots,\mathbbm{x}_i,\mathbbm{w}_0,\dots,\mathbbm{w}_{i-1})$.
    \item Computes $\pi_C^{\langle i \rangle} \gets g_1^{c_i(\tau)}$ with $c_i(X) := \delta_i^{-1}\sum_{j \in \mathcal{I}_W^{\langle i \rangle}}z_jf_j(X) + r_i\delta_d$.
    \item Updates accumulator as $a_{i+1} \gets \mathcal{H}(a_i,\pi_C^{\langle i \rangle})$.
    \item If $i<d$, for each $j \in \mathcal{I}_X^{\langle i+1 \rangle}$ sets $z_j \gets \mathcal{H}(a_{i+1},g_1^j)$.
\end{itemize}
After finishing $d$ rounds, the prover does the following:
\begin{itemize}
    \item Computes polynomial $h(X)$ from QAP as in Groth16.
    \item Samples random $r,s \gets \mathbb{F}$ and computes $\pi_A \gets g_1^{a(\tau)}$, $\pi_B \gets g_2^{b(\tau)}$ and the last proof piece $\pi_C^{\langle d \rangle} \gets g^{c_d(\tau)}$ where $a(X)$, $b(X)$, and $c(X)$ are defined as follows:
    \begin{align*}
        a(X) &= \alpha + \sum_{i \in [n]} z_i\ell_i(X) + r\delta_d, \\
        b(X) &= \beta + \sum_{i \in [n]}z_ir_i(X) + s\delta_d
    \end{align*}
    \[
    c_d(X) = \delta_d^{-1}\left(\sum_{i \in \mathcal{I}_W^{\langle d \rangle}}z_if_i(X) + h(X)t_{\Omega}(X)\right) + a(X)s+b(X)r-\sum_{i \in [d]}r_i\delta_i - rs\delta_d
    \]
    \item \textbf{Outputs} proof $\pi=(\pi_A,\pi_B,\pi_C^{\langle 0 \rangle}, \dots, \pi_C^{\langle d \rangle}) \in \mathbb{G}_1 \times \mathbb{G}_2 \times \mathbb{G}_1^{d+1}$.
\end{itemize}

\textbf{Verify($\mathsf{vp},\mathbbm{x}_0,\pi$)}. The verifier recomputes all the challenges $\{\mathbbm{x}_i\}_{i \in [d+1] \setminus \{0\}}$ and verifies that corresponding public signals are correct. Then, she checks:
\[
e(\pi_A,\pi_B) = e(g_1^{\alpha},g_2^{\beta}) \cdot e(g_1^{i(\tau)},g_2^{\gamma}) \cdot \prod_{i \in [d+1]}e(\pi_C^{\langle i \rangle}, g_2^{\delta_i})
\]

\subsection{Efficiency Analysis}

Finally, we specify the efficiency of UltraGroth.
\begin{itemize}
    \item The prover's work is essentially the same as in Groth16: if the total witness size is $n$, the prover needs to conduct $\mathcal{O}(n \cdot E_1 + n \cdot E_2)$ operations with a couple of additional hashes (formally, $\sum_{i \in [d+1] \setminus \{0\}}|\mathbbm{x}_i|$ hashes, which is a negligible amount in practice). Note however, that due to such R1CS splitting, the circuit's size $n$ is drastically decreased: in many cases it becomes sublinear in the original size.
    \item Verifier's work is $(d+3)E+\mathcal{O}(l \cdot E_1)$.
    \item Proof size is $1\;\mathbb{G}_2$ point and $(d+2)\;\mathbb{G}_1$ points.
\end{itemize}

\subsection{Practical Implementation and Further Steps}

Currently, we have implemented the UltraGroth as a part of the Bionetta project
for $d=1$ in the folowing repository:
\url{https://github.com/rarimo/ultragroth}. Also, we have a custom
\href{https://github.com/rarimo/ultragroth-snarkjs}{\textit{SnarkJS} fork} that
implements smart contract generation and witness verify and export operations.
But the current implementation is vastly inflexible and currently primarily
works in conjuction with \textit{Bionetta's witness generator} only. This
implementation takes the regular Circom-generated \texttt{.wtns} and
\texttt{.sym} files, parses them, figures out which witness element corresponds
to the randomness, and during the proving procedure fills in this spot using the
protocol above. 

We look forward to completing the implementation that could be easily integrated
into other's projects. One nice way to implement that would be to add a special
keyword to the Circom compiler. For instance, it would be nice if we could write
\texttt{random signal r[N]} that fills in the vector \texttt{r} with random
values, although such task is surely challenging.

\subsection{Acknowledgements}

Of course, without Lev Soukhanov this section would not exist. I also thank
Artem Sdobnov, Vitalii Volovyk, Yevhenii Sekhin, and Illia Dovgopoly for the
UltraGroth implementation for Bionetta.

\subsection{Summary}

In case you implement complex circuits involving non-native operations (e.g.,
wrapping the verification of some ZK protocol), it is definitely worth checking
out the UltraGroth protocol. With the tiny verification overhead, one gets an
immense boost to the proving procedure. In particular, it makes client-side ZK
much more practical and realistic without sacrificing verification efficiency.
